\documentclass[11pt]{article}
\usepackage{float}
\usepackage{amssymb, amsmath, amsfonts}
\begin{document}
\section{Students Exams Score Analysis }
You are preparing a data analysis report for a classroom dataset. The goal is to analyze student performance using vectors and basic statistics, then present it professionally in LaTeX. This aligns directly with data science practice: organizing data, computing summaries, and presenting results clearly.\\[6pt]
Dataset:
\begin{table} [H]
\centering
\def \arraystrech{3inc}
\begin{tabular}{|l|c|c|c|}
\hline
Student & Math &Phsic & Chem\\
\hline
Alice & 90 & 85 & 92\\
Bob & 78 & 82 & 88\\
Charlie & 85 &89 & 84\\
Diana & 92 & 95 & 90\\
Ethan & 70 & 75 &80\\ \hline 
\end{tabular}
\caption{Dataset table}
\end{table}
This is the computed and converted table:
\begin{table}[H]
\centering
\def \arraystrech{3cm}
\begin{tabular}{l || c c c}
\hline
Students & Marks Vector & Total $T_i$ & Average $m_i$\\ \hline
 Alice & $M_{Alice}\left[90, 85, 92\right]^T$ & 267 & 89.0 \\ \hline
 Bob & $M_{Bob} \left[78, 82, 88\right]^T$ & 248& 82.67\\ \hline
Charlie & $M_{Chalie} \left[85, 89, 84\right]^T$ &258& 86\\ \hline
Diana & $M_{Diana}\left[92, 95, 90\right]^T$ & 277 & 92.33\\ \hline
Ethan & $M_{Ethan}\left[70, 75 80\right]^T$ & 225 & 75\\ \hline
\end{tabular}
\end{table}
Finally, standard deviation for each student is computed by the following formula:
$$\sigma_i = \sqrt{\frac{1}{3}\sum_{j=1}^3 \left(m_{ij} - m_i\right)^2}$$
\end{document}